%% Compilar: latexmk -pdf apresentacao.tex  (pdflatex + biber).
\documentclass[
  aspectratio = 169,
  10pt,
  t,
  onlytextwidth,
  compress,
  handout,
  xcolor = table,
  english,
  brazilian,
]{beamer}

\usepackage[
  ABNTCit  = NSB,
  DocLic   = None,
  URLIcon  = On,
  DOIIcon  = On,
]{utfpr-slides}

\SlidesBackGround{utfpr-b-bg-16-x-9}
\SlidesLogo{utfpr-pg-d}

\newcommand{\cnxParams}{28,01}\newcommand{\cnxSize}{112,1}\newcommand{\rnParams}{24,03}\newcommand{\effParams}{4,34}
\newcommand{\cnxVal}{3,84}\newcommand{\cnxA}{7,87}\newcommand{\cnxC}{8,32}\newcommand{\cnxO}{10,42}
\newcommand{\rnVal}{5,36}\newcommand{\rnA}{9,54}\newcommand{\rnC}{13,20}\newcommand{\rnO}{14,78}
\newcommand{\effVal}{5,10}\newcommand{\effA}{7,92}\newcommand{\effC}{12,24}\newcommand{\effO}{14,08}
\newcommand{\ablVal}{4,38}\newcommand{\ablA}{6,89}\newcommand{\ablC}{10,35}\newcommand{\ablO}{11,29}
\newcommand{\imVal}{27,50}\newcommand{\imA}{29,49}\newcommand{\imC}{34,88}\newcommand{\imO}{34,75}
\newcommand{\fusVal}{1,49}\newcommand{\fusA}{2,73}\newcommand{\fusC}{4,96}\newcommand{\fusO}{6,86}
\newcommand{\cnxCIA}{5,6 a 10,0\% (FVC2002)}\newcommand{\cnxCIC}{6,3 a 11,0\% (FVC2004)}
\newcommand{\cnxFMRA}{6,4\%}\newcommand{\cnxFMRC}{6,9\%}\newcommand{\cnxFNMRA}{9,6\%}\newcommand{\cnxFNMRC}{10,8\%}
\newcommand{\cnxThr}{0,368}\newcommand{\cnxThrFig}{0,380}
\newcommand{\custo}{Custo por imagem: rede em 9,7 / 10,3 / 16,3~ms (ConvNeXt / ResNet / EfficientNet, GPU, lote 1, média); pré-processamento com realce de Gabor, 76 a 125~ms em CPU.}
\newcommand{\cnxTransfer}{EER de 3,8\% na validação, 7,9\% no teste da FVC2002 e 8,3\% na FVC2004: a perda maior é dentro da própria base. Na FVC2004 o erro dominante é de rejeição (FNMR de 10,8\%).}
\newcommand{\robBlur}{Desfoque com $\sigma=3$: EER de 27,2\% e \alert{FMR de 75,7\%}; o sistema aceita impostores.}
\newcommand{\robRot}{Rotação de $30^\circ$: EER de 10,7\%; só a FNMR sobe (18,5\%).}
\newcommand{\ablText}{Sem realce, a EER cai 1,0 ponto na FVC2002 (IC contém zero) e sobe 2,0 pontos na FVC2004 (IC de 0,3 a 3,4). No protocolo oficial o ganho é de 0,9 ponto e o IC contém zero.}
\newcommand{\errFM}{Falsos matches com score de 0,78 a 0,81, bem acima do limiar de 0,380: laços com núcleo e fluxo de cristas parecidos (dedos 92 e 98, 23 e 36).}
\newcommand{\errFNM}{Os três falsos não-matches envolvem 75\_5, que mostra só a ponta do dedo (99~mm$^2$) e fica sem região comum com as outras impressões.}
\newcommand{\errC}{Na FVC2004 DB1 os erros extremos repetem o quadro: capturas parciais e imagens pouco nítidas.}
\newcommand{\conclA}{Na FVC2002 DB1 as três CNNs empatam; a ConvNeXt-Tiny tem EER de 7,9\% (IC de 5,6 a 10,0).}
\newcommand{\conclB}{Na FVC2004 DB1, nunca treinada, a ConvNeXt-Tiny fica de 3,7 a 4,9 pontos abaixo das outras, com o mesmo limiar.}
\newcommand{\conclC}{O SourceAFIS ganha na FVC2002 (2,7\%), mas a fusão com a CNN tem a menor EER da configuração base na FVC2004 (5,0\% e 6,9\%) e iguala o SourceAFIS na FVC2002 (2,7\%).}
\newcommand{\conclD}{Com desfoque forte o sistema aceita 76\% dos impostores: o QC precisa de um critério de nitidez.}


\title[Impressão digital 1:1 com ConvNeXt-Tiny]{%
  Verificação 1:1 de impressões digitais\\com embeddings de CNN%
}
\subtitle{ConvNeXt-Tiny, ResNet-50 e EfficientNet-B0 na FVC2002 DB1 e na FVC2004 DB1}
\subject{Projeto Final: Biometria (UTFPR Ponta Grossa)}
\author[L.~R.~de Carvalho]{%
  Lucas Ribeiro de Carvalho%
  \AuthorInfo{%
    ORCID = {0009-0005-2339-0093},%
    Affil = {UTFPR},%
    Inst  = {1},%
  }%
  \Inst{\faUserTalk}%
}
\institute[UTFPR]{%
  \Affil[1]{\UTFPRName, Ponta Grossa, Paraná, Brasil}%
  \and\InfoList[\tiny]{\faOrcid}{\MakeInfoList{ORCID}}%
}
\DateOrInfo[Biometria, \YearNum]{Disciplina de Biometria, Dupla 3, \YearNum}
\titlegraphic{%
  \hfill%
  \Logo{public-education-l}\hspace*{2mm}%
  \Logo{utfpr-pg-l}%
}
\Campus{Campus Ponta Grossa}
\Department[Biometria]{Disciplina de Biometria}

\addbibresource{referencias.bib}

\newcommand{\fig}[2]{\includegraphics[#1]{../artigo/figures/#2}}

\begin{document}

\maketitle

\section{Problema}

\begin{frame}{Problema e pergunta}
\begin{columns}[T]
\column{0.55\textwidth}
\begin{itemize}
\item Verificação 1:1: duas impressões, mesmo dedo ou não?
\item A CNN gera um \alert{embedding de 256 dimensões}; a comparação é a similaridade de cosseno contra um limiar.
\item Modelos de tamanho fixo publicados usam centenas de milhares de imagens~\cite{engelsma2021deepprint}.
\item Aqui o ajuste fino tem \alert{480 imagens de 60 dedos}.
\end{itemize}
\column{0.42\textwidth}
\begin{block}{Pergunta}
Como três CNNs pré-treinadas na ImageNet se comportam com pouco dado, quanto do desempenho sobrevive à troca de base e sensor, e quanto custam?
\end{block}
\end{columns}
\end{frame}

\section{Método}

\begin{frame}{Dados e protocolo}
\begin{columns}[T]
\column{0.55\textwidth}
\footnotesize
\begin{tabular}{@{}llc@{}}
\toprule
Papel & Base & Dedos \\
\midrule
Treino & FVC2002 DB1~\cite{maio2002fvc2002} (A+B) & 60 \\
Validação & FVC2002 DB1 & 20 \\
Teste A & FVC2002 DB1 & 30 \\
Teste C & FVC2004 DB1~\cite{maio2004fvc2004}, nunca treinada & 30 \\
Teste C, oficial & FVC2004 DB1-A & 100 \\
\bottomrule
\end{tabular}
\column{0.42\textwidth}
\begin{itemize}
\item Divisão por identidade, seed 42: nenhum dedo em duas partes.
\item Época e limiar só na validação; o \alert{limiar fica congelado}.
\item Treinos com as seeds 2026, 1999 e 1234.
\end{itemize}
\end{columns}
\par\medskip
\footnotesize
Sensores ópticos de 500~dpi: Identix TouchView~II (FVC2002 DB1) e CrossMatch V300 (FVC2004 DB1). 110 dedos $\times$ 8 impressões por base.\\
Pares: validação 560 genuínos e 12\,160 impostores; cada teste de 30 dedos 840 e 27\,840; protocolo oficial 2\,800 e 4\,950.
\end{frame}

\begin{frame}{Pipeline}
\centering
\fig{width=0.76\textwidth}{pipeline.png}
\par\medskip
\footnotesize
500~dpi $\rightarrow$ ROI por blocos $\rightarrow$ QC (FTA se ROI $<40$~mm$^2$) $\rightarrow$ normalização de Hong $\rightarrow$ realce de Gabor~\cite{hong1998enhancement} $\rightarrow$ janela $320\times320$ $\rightarrow$ CNN $\rightarrow$ embedding $\rightarrow$ cosseno $\rightarrow$ limiar $\rightarrow$ decisão
\par\smallskip
Desvios do pipeline do enunciado: o realce de Gabor (ablação nos resultados) e a janela fixa no lugar do redimensionamento, para manter a escala das cristas entre sensores.
\end{frame}

\begin{frame}{CNN, embedding e matching}
\begin{columns}[T]
\column{0.5\textwidth}
\begin{block}{ConvNeXt-Tiny~\cite{liu2022convnext} (principal)}
\begin{itemize}
\item ImageNet-1k (\texttt{IMAGENET1K\_V1}), 1 canal.
\item \texttt{features.0}--\texttt{3} congeladas.
\item \cnxParams~M parâmetros, \cnxSize~MB.
\end{itemize}
\end{block}
\begin{block}{Comparação}
\begin{itemize}
\item ResNet-50~\cite{he2016resnet}: \texttt{layer1}--\texttt{2} congeladas, \rnParams~M, 96,4~MB.
\item EfficientNet-B0~\cite{tan2019efficientnet}: \texttt{features.0}--\texttt{3} congeladas, \effParams~M, 17,6~MB.
\end{itemize}
\end{block}
\column{0.46\textwidth}
\begin{itemize}
\item Cabeça: Linear $\rightarrow$ 256 + BatchNorm, norma L2.
\item Perda ArcFace~\cite{deng2019arcface}, $m=0{,}3$, $s=30$; lotes de 16 dedos $\times$ 4.
\item AdamW, lr $3\times10^{-4}$, 80 épocas; melhor época pela EER de validação.
\item Score: $s=\mathbf{e}_a^\top\mathbf{e}_b$; aceita se $s\ge t$.
\item Hardware: RTX~4060 8~GB, Ryzen~5 5600X, PyTorch 2.6; um treino leva de 3,4 a 5,7~min.
\end{itemize}
\end{columns}
\end{frame}

\section{Resultados}

\begin{frame}{Resultados: três CNNs (A, B e C)}
\begin{columns}[T]
\column{0.5\textwidth}
\scriptsize
\setlength{\tabcolsep}{2.5pt}
\begin{tabular}{@{}lccc@{}}
\toprule
EER (\%) & FVC2002 & FVC2004 & 2004 oficial \\
\midrule
\alert{ConvNeXt-Tiny} & \alert{$7{,}87\pm0{,}29$} & \alert{$8{,}32\pm0{,}44$} & \alert{$10{,}42\pm0{,}17$} \\
ResNet-50 & $9{,}54\pm0{,}95$ & $13{,}20\pm0{,}64$ & $14{,}78\pm1{,}02$ \\
EfficientNet-B0 & $7{,}92\pm0{,}40$ & $12{,}24\pm0{,}39$ & $14{,}08\pm0{,}46$ \\
ConvNeXt s/ realce & $6{,}89\pm0{,}78$ & $10{,}35\pm0{,}72$ & $11{,}29\pm0{,}84$ \\
ConvNeXt s/ ajuste & \imA & \imC & \imO \\
\bottomrule
\end{tabular}
\par\medskip
\footnotesize
Média $\pm$ desvio-padrão de três seeds. EER de validação: \cnxVal, \rnVal\ e \effVal.\\
IC 95\% da ConvNeXt (bootstrap sobre dedos): \cnxCIA\ e \cnxCIC.\\
Experimento B parcial: sem a CNN da outra dupla.
\column{0.47\textwidth}
\scriptsize
\setlength{\tabcolsep}{2.5pt}
\begin{tabular}{@{}lccc@{}}
\toprule
ConvNeXt-Tiny & FVC2002 & FVC2004 & 2004 oficial \\
\midrule
\alert{Limiar} & \multicolumn{3}{c}{\alert{$0{,}368\pm0{,}010$} (congelado)} \\
FMR (\%) & 6,39 & 6,92 & 7,57 \\
FNMR (\%) & 9,56 & 10,79 & 14,10 \\
AUC & 0,976 & 0,966 & 0,955 \\
TAR@FAR=1\% (\%) & 68,8 & 46,5 & 50,9 \\
FTA (\%) & 0,00 & 0,00 & 0,00 \\
\bottomrule
\end{tabular}
\par\medskip
\footnotesize
Rede, GPU, lote 1 (média): 9,7 / 10,3 / 16,3~ms para ConvNeXt / ResNet / EfficientNet.\\
Pré-processamento com realce: 76 a 125~ms em CPU.
\end{columns}
\end{frame}

\begin{frame}{Resultados: minúcias e fusão}
\begin{columns}[T]
\column{0.54\textwidth}
\scriptsize
\setlength{\tabcolsep}{3pt}
\begin{tabular}{@{}lcccc@{}}
\toprule
EER (\%) & Valid. & FVC2002 & FVC2004 & 2004 oficial \\
\midrule
ConvNeXt-Tiny & \cnxVal & \cnxA & \cnxC & \cnxO \\
SourceAFIS~\cite{vazan2024sourceafis} & 2,61 & 2,71 & 9,31 & 12,18 \\
\alert{Fusão CNN + minúcias} & \fusVal & \alert{\fusA} & \alert{\fusC} & \alert{\fusO} \\
\bottomrule
\end{tabular}
\par\medskip
\footnotesize
\begin{itemize}
\item O SourceAFIS ganha com folga na FVC2002 DB1 e é quem mais piora na troca de sensor.
\item Fusão: scores normalizados pelos impostores de validação, peso de 0,6 ou 0,7 para a CNN, fixado na validação.
\item Na FVC2004 DB1 a fusão fica abaixo dos dois comparadores, com intervalos que excluem zero.
\end{itemize}
\column{0.42\textwidth}
\centering
\fig{width=\linewidth}{roc.pdf}\\[-2pt]
{\scriptsize ROC da seed 2026}
\end{columns}
\end{frame}

\begin{frame}{Scores e limiar}
\centering
\fig{width=0.8\textwidth}{distributions.pdf}
\par\medskip
\begin{itemize}
\item Limiar da validação reutilizado sem mudança nas duas bases: \cnxThr\ na média das três seeds; a figura é da seed 2026 (\cnxThrFig).
\item \cnxTransfer
\end{itemize}
\end{frame}

\begin{frame}{Robustez e ablação}
\begin{columns}[T]
\column{0.42\textwidth}
\centering
\fig{height=0.72\textheight}{robustness.pdf}
\column{0.54\textwidth}
\begin{itemize}
\item Desfoque e rotação em quatro níveis nos 30 dedos de teste da FVC2002 DB1, modelo e limiar congelados.
\item \robBlur
\item \robRot
\item O QC por área marcou só 3 das 240 imagens com $\sigma=3$.
\end{itemize}
\begin{block}{Ablação: sem o realce de Gabor}
\ablText
\end{block}
\end{columns}
\end{frame}

\begin{frame}{Análise de erros}
\centering
\fig{width=0.66\textwidth}{errors_fvc2002_db1.png}\\[-2pt]
{\scriptsize FVC2002 DB1: falsos matches (acima) e falsos não-matches (abaixo)}
\par\smallskip
\footnotesize
\begin{itemize}
\item \errFM
\item \errFNM
\item \errC
\end{itemize}
\end{frame}

\section{Limitações}

\begin{frame}{Limitações}
\begin{itemize}
\item \alert{Experimento B parcial}: os scores da outra dupla não chegaram; a comparação é entre as nossas três CNNs.
\item Treino com 60 dedos de um só sensor; o limiar e a época saem de 20 dedos, e a EER de validação (3,8\%) é metade da de teste.
\item Os hiperparâmetros vêm de uma fase anterior em que os dedos 107 e 109 estavam no teste. Sem eles, a EER é 7,95\% na FVC2002 DB1 e 7,94\% na FVC2004 DB1.
\item A ResNet-50 e a EfficientNet-B0 não tiveram ajuste próprio; a latência do SourceAFIS não foi medida.
\item Desfoque e rotação medidos só na FVC2002 DB1; as hipóteses da análise de erros não foram testadas.
\end{itemize}
\end{frame}

\section{Conclusões}

\begin{frame}{Conclusões}
\begin{itemize}
\item \conclA
\item \conclB
\item \conclC
\item \conclD
\end{itemize}
\vfill
\begin{exampleblock}{Próximos passos}
\small Mais dados de treino (outras bases FVC), critério de nitidez no QC, período de crista por região e alinhamento por pose.
\end{exampleblock}
\par\smallskip
{\footnotesize Código, protocolos e scores: \url{https://github.com/lucthesloth/biometria26.2}, tag \texttt{final-v1.0}.}
\end{frame}

\section<presentation>{\refname}

\begin{frame}<presentation>[allowframebreaks]{\refname}
\printbibliography[heading = none]
\end{frame}

\appendix

\begin{frame}{Apoio: otimização}
\begin{columns}[T]
\column{0.5\textwidth}
\scriptsize
\setlength{\tabcolsep}{3pt}
\begin{tabular}{@{}lcc@{}}
\toprule
EER / FMR / FNMR (\%) & FVC2002 & 2004 oficial \\
\midrule
Base: 60 dedos, 1 rede & 7,9 / 6,4 / 9,6 & 10,4 / 7,6 / 14,1 \\
390 dedos, 2 redes & 4,1 / 3,6 / 4,6 & 6,2 / 3,5 / 11,0 \\
745 dedos, 4 redes & 3,1 / 3,2 / 3,1 & 4,3 / 2,3 / 8,1 \\
\alert{745 dedos, 12 redes} & \alert{2,7 / 2,8 / 2,7} & \alert{3,8 / 1,9 / 7,5} \\
\bottomrule
\end{tabular}
\par\medskip
\footnotesize
\begin{itemize}
\item À parte do resultado principal, só para a ConvNeXt-Tiny.
\item Ficaram: mais bases no treino (a \alert{FVC2004 DB1 segue fora}), soma de redes com Gabor e com Hong, score normalizado.
\item Saíram 10 tentativas que não melhoraram a validação.
\item Custo: 12 redes, cerca de 110~ms e 1,3~GB.
\end{itemize}
\column{0.46\textwidth}
\centering
\fig{height=0.74\textheight}{otimizacao.pdf}
\end{columns}
\end{frame}

\begin{frame}{Apoio: tentativas}
\centering
\scriptsize
\setlength{\tabcolsep}{4pt}
\begin{tabular}{@{}lccl@{\hspace{8mm}}lccl@{}}
\toprule
Etapa 1 & Sem & Com & & Etapa 2 & Sem & Com & \\
\midrule
Apagamento aleatório & 3,41 & 3,04 & fora & Distorção suave, cristas finas/grossas & 2,70 & 3,45 & fora \\
Apagamento e translação 20\% & 3,41 & 2,60 & fora & Média móvel dos pesos & 2,70 & 3,04 & fora \\
Um estágio congelado a menos & 3,41 & 3,21 & fora & Busca de rotação na comparação & 2,27 & 2,32 & fora \\
Margem ArcFace 0,5 & 3,41 & 3,75 & fora & 785 dedos (FVC2004 DB2--4) & 2,70 & 2,20 & \alert{fica} \\
160 épocas & 3,41 & 3,29 & fora & Rotação $45^\circ$, translação 20\% & 5,78 & 5,18 & \alert{fica} \\
390 dedos & 3,84 & 2,27 & \alert{fica} & Translação 30\% & 5,18 & 5,71 & fora \\
Soma Gabor + Hong & 2,27 & 1,86 & \alert{fica} & Score normalizado & 3,39 & 3,10 & \alert{fica} \\
Média em rotações & 2,70 & 3,21 & fora & Soma das três seeds & 2,67 & 2,26 & \alert{fica} \\
\bottomrule
\end{tabular}
\par\medskip
\footnotesize
EER (\%) nos dedos de validação, sem e com cada mudança. Na etapa 1, ganhos abaixo de um ponto com uma seed ficaram fora.
\end{frame}

\end{document}
